% Generated draft: requires mathematical and source review.
\documentclass[11pt]{article}
\usepackage[margin=25mm]{geometry}
\usepackage{fontspec}
\setmainfont{Latin Modern Roman}
\usepackage{amsmath,amssymb,mathtools,unicode-math}
\setmathfont{Latin Modern Math}
\usepackage{enumitem,needspace,xurl,hyperref}
\setlength{\parindent}{0pt}
\setlength{\parskip}{5pt}
\setlength{\emergencystretch}{3em}
\newcommand{\N}{\mathbb{N}}
\newcommand{\Z}{\mathbb{Z}}
\newcommand{\Q}{\mathbb{Q}}
\newcommand{\R}{\mathbb{R}}
\newcommand{\C}{\mathbb{C}}
\newcommand{\F}{\mathbb{F}}
\newcommand{\A}{\mathbb{A}}
\newcommand{\PP}{\mathbb{P}}
\newcommand{\E}{\mathbb{E}}
\newcommand{\eps}{\varepsilon}
\newcommand{\dd}{\,\mathrm{d}}
\begin{document}
\section*{UnsolvedMath ID: 701454}
\textbf{Source problem code:} catalogue entry
\subsection{Definitions and mathematical statement}
All groups and formations below are understood in the finite-group setting. A formation is a nonempty class of finite groups, closed under isomorphism, homomorphic images, and finite subdirect products. Here a subdirect product of groups \(H_1,\ldots,H_n\) is a subgroup of \(H_1\times\cdots\times H_n\) whose projection onto each \(H_i\) is surjective. Write \(\mathrm{form}(G)\) for the smallest formation containing \(G\). A subformation of \(\mathrm{form}(G)\) is a formation contained in \(\mathrm{form}(G)\). Write \(Z_r\) for a cyclic group of prime order \(r\).

A composition series is a strictly increasing chain from the trivial subgroup to the whole group, with each term normal in the next and each consecutive quotient simple. Let \(G\) be a finite group having exactly one composition series, namely
\[
1\triangleleft G_1\triangleleft G_2\triangleleft G,
\]
and suppose that
\[
G_1\cong Z_p,\qquad G_2/G_1\cong X,\qquad G/G_2\cong Z_q,
\]
where \(p\) and \(q\) are primes and \(X\) is a nonabelian finite simple group. Uniqueness here means uniqueness of the chain of subgroups, not merely uniqueness of the composition factors up to isomorphism and order.

Is it true that, for every such \(G\),
\[
\mathrm{form}(G)\text{ has infinitely many distinct subformations}
\quad\Longleftrightarrow\quad p=q\,?
\]
\subsection{Short English statement}
For a finite group with a unique composition series whose factors, in order from bottom to top, are cyclic of order \(p\), nonabelian simple, and cyclic of order \(q\), ask whether its generated formation has infinitely many subformations exactly when \(p=q\).
\subsection{Sources}
\begin{itemize}
\item Local UnsolvedMath registry, supplied complete record: permanent ID \(701454\), exact source code \texttt{catalogue entry}, title \emph{Kourovka Notebook Problem 18.25}; upstream revision \texttt{d371ee917ac9647fbb06cd6fc529deec6f3dda82}.
\item The supplied record attributes Problem 18.25 to V. P. Burichenko, proposed in 2014, and cites \emph{The Kourovka Notebook (Unsolved Problems in Group Theory)}, 18th issue (2014), and the 21st issue edition (2026), \texttt{arXiv:1401.0300v48}. The collection description names E. I. Khukhro and V. D. Mazurov as editors.
\item The supplied registry labels the problem \emph{open} and records an update dated 2026-10-16. This is supplied source metadata, not an independently verified assessment of the literature or mathematical status.
\end{itemize}
\end{document}
